\documentclass[11pt]{article}

\usepackage[T1]{fontenc}
\usepackage[utf8]{inputenc}
\usepackage{lmodern}
\usepackage[margin=1in]{geometry}
\usepackage{amsmath,amssymb,mathtools}
\usepackage{booktabs}
\usepackage{xcolor}
\usepackage[
    colorlinks=true,
    linkcolor=blue,
    urlcolor=blue,
    citecolor=blue
]{hyperref}

\title{Theory Notes for the 1D Silicon p--n Junction Solver}
\author{}
\date{}

\begin{document}

\maketitle
\tableofcontents
\bigskip

\section{Overview}

This document describes the mathematical model and numerical procedure used in
the DOLFINx-based one-dimensional p--n junction solver in the repository.
The code solves a steady-state drift-diffusion--Poisson system on a one-dimensional
interval, using a mixed finite-element formulation for the electrostatic
potential and the two quasi-Fermi energies.

The current implementation is intentionally simple:
\begin{itemize}
    \item the device is one-dimensional in space,
    \item the doping profile is a smooth p-to-n transition,
    \item Maxwell--Boltzmann statistics are used,
    \item the equations are written internally in atomic units,
    \item each bias point is solved by Newton iteration, using the previous bias
    point as the initial guess.
\end{itemize}

\section{Unknowns and Domain}

The spatial domain is the interval
\[
    x \in [0,L].
\]

The unknown fields are
\[
    \phi(x), \qquad F_n(x), \qquad F_p(x),
\]
where
\begin{itemize}
    \item $\phi(x)$ is the electrostatic potential,
    \item $F_n(x)$ is the electron quasi-Fermi energy,
    \item $F_p(x)$ is the hole quasi-Fermi energy.
\end{itemize}

The code uses atomic units internally, so the energy-like quantities
$\phi$, $F_n$, and $F_p$ are represented in Hartree units.

\section{Unit System}

\subsection{Atomic units used internally}

The solver uses:
\begin{align*}
    a_0 &= 5.29177210903 \times 10^{-11}\ \mathrm{m}, \\
    E_h &= 27.211386245988\ \mathrm{eV}, \\
    t_0 &= 2.4188843265857 \times 10^{-17}\ \mathrm{s}, \\
    k_B &= 3.166811563 \times 10^{-6}\ E_h/\mathrm{K}.
\end{align*}

At $T=300\ \mathrm{K}$,
\[
    k_B T \approx 9.500434689 \times 10^{-4}\ E_h.
\]

\subsection{Derived unit conventions}

The code represents:
\begin{itemize}
    \item length in units of $a_0$,
    \item energy and electrostatic potential in units of $E_h$,
    \item time in units of $t_0$,
    \item number density in units of $a_0^{-3}$.
\end{itemize}

The plots are converted back to more readable units:
\begin{itemize}
    \item position in microns,
    \item doping in $\mathrm{m}^{-3}$,
    \item applied bias in volts,
    \item current in amperes.
\end{itemize}

\section{Carrier Statistics}

The solver uses Maxwell--Boltzmann carrier densities:
\begin{align}
    n(\phi,F_n) &= N_c \exp\!\left(\frac{F_n + \phi - E_{c0}}{k_B T}\right), \\
    p(\phi,F_p) &= N_v \exp\!\left(\frac{E_{v0} - \phi - F_p}{k_B T}\right).
\end{align}

To control overflow in the numerical implementation, the exponent arguments are
clipped to the interval $[-80,80]$ before exponentiation.

\section{Doping Model}

The device is a p--n junction with a smooth transition centered at $x=L/2$.
The donor and acceptor profiles are
\begin{align}
    N_D(x) &= \frac{N_D^\star}{2}
    \left[
        1 + \tanh\!\left(\frac{x-L/2}{w}\right)
    \right], \\
    N_A(x) &= \frac{N_A^\star}{2}
    \left[
        1 - \tanh\!\left(\frac{x-L/2}{w}\right)
    \right].
\end{align}

This produces a p-type left side and an n-type right side while avoiding a
discontinuous step in the finite-element interpolation.

\section{Governing Equations}

\subsection{Poisson equation}

The electrostatic equation is written as
\begin{equation}
    -\frac{d}{dx}
    \left(
        \epsilon_r \frac{d\phi}{dx}
    \right)
    =
    A \bigl(p - n + N_D - N_A\bigr).
\end{equation}

Here $A$ is the effective device cross-sectional area used in the 1D reduction.
The code keeps the densities as volumetric densities and multiplies the source
term by the cross-sectional area when reducing the device to one dimension.

\subsection{Continuity equations}

The drift-diffusion currents are modeled by
\begin{align}
    J_n &= \mu_n\, n \frac{dF_n}{dx}, \\
    J_p &= \mu_p\, p \frac{dF_p}{dx}.
\end{align}

The steady-state continuity equations are represented in weak form, but the
strong-form sign structure corresponds to
\begin{align}
    \frac{dJ_n}{dx} &= -A\,U, \\
    \frac{dJ_p}{dx} &= +A\,U,
\end{align}
where $U$ is the recombination source term defined below.

\subsection{Recombination model}

The code uses the simple recombination model
\begin{equation}
    U = \gamma \bigl(np - n_i^2\bigr).
\end{equation}

To stay consistent with the band-edge parameters used in the code, the intrinsic
density entering this term is the mass-action value
\begin{equation}
    n_i = \sqrt{N_c N_v}
    \exp\!\left(
        \frac{E_{v0}-E_{c0}}{2k_B T}
    \right).
\end{equation}

\section{Variational Formulation}

Let $v_\phi$, $v_n$, and $v_p$ be the test functions associated with
$\phi$, $F_n$, and $F_p$, respectively.

\subsection{Poisson weak form}

The weak form used in the code is
\begin{equation}
    \int_0^L
    \epsilon_r
    \frac{d\phi}{dx}
    \frac{dv_\phi}{dx}\,dx
    -
    \int_0^L
    A\,(p-n+N_D-N_A)\,v_\phi\,dx
    = 0.
\end{equation}

\subsection{Electron continuity weak form}

The electron residual is
\begin{equation}
    \int_0^L
    \mu_n\,n\,\frac{dF_n}{dx}\frac{dv_n}{dx}\,dx
    +
    \int_0^L
    A\,U\,v_n\,dx
    = 0.
\end{equation}

\subsection{Hole continuity weak form}

The hole residual is
\begin{equation}
    -\int_0^L
    \mu_p\,p\,\frac{dF_p}{dx}\frac{dv_p}{dx}\,dx
    +
    \int_0^L
    A\,U\,v_p\,dx
    = 0.
\end{equation}

\subsection{Mixed nonlinear system}

The total nonlinear residual is the sum of the three weak forms:
\[
    \mathcal{R}(\phi,F_n,F_p;v_\phi,v_n,v_p)
    =
    \mathcal{R}_\phi + \mathcal{R}_n + \mathcal{R}_p.
\]

The code solves this coupled residual directly using DOLFINx and PETSc.

\section{Boundary Conditions}

\subsection{Ohmic contact model}

The contacts are modeled as ideal Ohmic Dirichlet boundaries for all three
unknowns:
\[
    \phi(0),\ \phi(L),\ F_n(0),\ F_n(L),\ F_p(0),\ F_p(L).
\]

\subsection{Equilibrium contact values}

At the left p-contact the code assumes approximate charge neutrality:
\begin{align}
    p_L &\approx N_A^\star, \\
    n_L &\approx \frac{n_i^2}{N_A^\star}.
\end{align}

At the right n-contact:
\begin{align}
    n_R &\approx N_D^\star, \\
    p_R &\approx \frac{n_i^2}{N_D^\star}.
\end{align}

The quasi-Fermi reference level is built by inverting the Maxwell--Boltzmann
relations at the contact:
\begin{align}
    F_n &= E_{c0} - \phi + k_B T \ln\!\left(\frac{n}{N_c}\right), \\
    F_p &= E_{v0} - \phi - k_B T \ln\!\left(\frac{p}{N_v}\right).
\end{align}

The code combines the electron-based and hole-based contact electrostatic
potentials by averaging the two resulting expressions at each contact.

\subsection{Applied bias}

For an applied voltage $V$, the current implementation uses the convention
\begin{align}
    \phi_L &= \phi_{L,\mathrm{eq}}, \\
    \phi_R &= \phi_{R,\mathrm{eq}} + V/E_h, \\
    F_{n,L} &= F_0, \\
    F_{p,L} &= F_0, \\
    F_{n,R} &= F_0 - V/E_h, \\
    F_{p,R} &= F_0 - V/E_h.
\end{align}

In other words, the left contact stays fixed at equilibrium, while the right
contact receives the bias shift.

\section{Constants Used in the Current Solver}

\subsection{Material and model constants}

\begin{center}
\begin{tabular}{lll}
\toprule
Quantity & Value in code & Meaning \\
\midrule
$\epsilon_r$ & $11.7$ & silicon relative permittivity \\
$E_{c0}$ & $1.12/27.211386245988\ E_h$ & conduction-band reference \\
$E_{v0}$ & $0$ & valence-band reference \\
$N_c$ & $2.8 \times 10^{25}\ \mathrm{m}^{-3}$ & conduction-band DOS \\
$N_v$ & $1.04 \times 10^{25}\ \mathrm{m}^{-3}$ & valence-band DOS \\
$n_i$ & $1.45 \times 10^{16}\ \mathrm{m}^{-3}$ & reported intrinsic density \\
$N_A^\star$ & $1.0 \times 10^{23}\ \mathrm{m}^{-3}$ & left-side doping scale \\
$N_D^\star$ & $1.0 \times 10^{23}\ \mathrm{m}^{-3}$ & right-side doping scale \\
$\gamma$ & $1.0 \times 10^{-20}\ \mathrm{m}^{3}/\mathrm{s}$ & recombination parameter \\
\bottomrule
\end{tabular}
\end{center}

\subsection{Geometry and mesh}

\begin{center}
\begin{tabular}{lll}
\toprule
Quantity & Value in code & Meaning \\
\midrule
$L$ & $2.0 \times 10^{-6}\ \mathrm{m}$ & device length \\
$A$ & $1.0\ \mu\mathrm{m}^2$ & effective cross-sectional area \\
$w$ & $2.0 \times 10^{-8}\ \mathrm{m}$ & junction transition width \\
cells & $200$ & 1D interval cells \\
\bottomrule
\end{tabular}
\end{center}

\subsection{Transport scaling}

The present code uses simplified mobility scaling rather than a detailed atomic-unit
conversion of the SI mobilities:
\begin{align}
    \mu_n &= 1.0, \\
    \mu_p &= 0.048/0.135.
\end{align}

This keeps the hole mobility proportional to the electron mobility, but the
values should be interpreted as pragmatic numerical scales, not as a full
microscopic transport calibration.

\section{Bias Sweep Procedure}

The applied biases used in the current code are
\[
    -0.5,\,-0.4,\,-0.3,\,-0.2,\,-0.1,\,
    0.0,\,
    0.1,\,0.2,\,0.3,\,0.4,\,0.5,\,0.6,\,0.7
    \quad \text{volts}.
\]

The actual loop is ordered by increasing magnitude of the applied bias:
\[
    0,\ -0.1,\ 0.1,\ -0.2,\ 0.2,\ \dots
\]
because continuation is easier when the next point is close to the previous one.

For each new bias point:
\begin{enumerate}
    \item the previous mixed solution is reused,
    \item the Dirichlet traces are shifted to the new contact values,
    \item Newton iteration is applied to the new residual,
    \item the resulting fields are plotted and the right-contact current is recorded.
\end{enumerate}

\section{Numerical Procedure}

\subsection{Finite-element spaces}

The code uses first-order Lagrange finite elements:
\[
    \phi,\ F_n,\ F_p \in \mathcal{V}_h,
\]
with a mixed product space
\[
    \mathcal{V}_h \times \mathcal{V}_h \times \mathcal{V}_h.
\]

\subsection{Nonlinear solver}

PETSc SNES is used through DOLFINx with:
\begin{itemize}
    \item Newton line search,
    \item backtracking line search,
    \item direct LU factorization in the linear subproblem,
    \item a small nonzero LU pivot shift for robustness.
\end{itemize}

The solver tolerances currently used are:
\begin{align*}
    \texttt{rtol} &= 10^{-8}, \\
    \texttt{atol} &= 10^{-10}, \\
    \texttt{max\_it} &= 80.
\end{align*}

\subsection{Initialization strategy}

At the first bias point, the code first solves the equilibrium Poisson-only
problem and uses that result to initialize $\phi$. The quasi-Fermi energies are
initialized by linear interpolation between the two contact values.

At later bias points, the full previous mixed solution is reused and shifted to
match the new boundary values.

\section{Outputs}

The current solver writes:
\begin{itemize}
    \item one bias plot per applied voltage,
    \item an $I$--$V$ curve,
    \item a current diagnostic plot.
\end{itemize}

Each bias plot contains:
\begin{itemize}
    \item $\phi(x)$,
    \item $F_n(x)$ and $F_p(x)$,
    \item $N_D(x)$ and $N_A(x)$.
\end{itemize}

\section{Current State of the Solver}

The equilibrium point is the most reliable part of the computation.
The biased drift-diffusion solve is still a numerically rough baseline rather
than a polished production-quality semiconductor solver.

In particular:
\begin{itemize}
    \item the continuation structure is in place,
    \item the equations and weak forms are explicit,
    \item the code is clean enough to iterate on,
    \item but the biased nonlinear solve still deserves further stabilization and
    physical calibration.
\end{itemize}

That is acceptable for a first educational and exploratory implementation, but
it should be treated as a starting point rather than a final validated device model.

\end{document}
